% GENERATED FILE. Edit canonical lessons, then run npm run generate:books.
% canonical-source-hash: fnv1a64:51c9de43576d3ea9
% canonical-lessons: FR-C166-photographe, FR-C166-journaliste, FR-C166-architecte, FR-C166-dentiste, FR-C166-pharmacien

\chapter{Photographer, Journalist, Architect, Dentist, Pharmacist}
\label{ch:fr-photographer-journalist-architect-dentist-pharmacist}
\hlchaptermodality{166}

\begin{chapteropening}
\textbf{By the end of this chapter:} \emph{I can say a photographer, a journalist, an architect, a dentist and a pharmacist.}

Complete the last lesson of chapter 166: I can say a photographer, a journalist, an architect, a dentist and a pharmacist.
\end{chapteropening}

\section[le photographe]{le photographe — a photographer}

\label{lesson:FR-C166-photographe}



\begin{quote}
Before the new one: say the French for a soldier, then the French for a tradition.
\end{quote}

\subsection*{You'll want to know: le photographe}
\textbf{le photographe} — \textquotedblleft{}a photographer\textquotedblright{}.

\begin{grammarlens}[title={What you've built}]
One more word for this chapter.
\end{grammarlens}

\subsection*{Your turn}
\begin{itemize}
\raggedright
  \item \emph{Say it:} \emph{le photographe}
  \item \emph{Say it:} \emph{le photographe}, once more
  \item \emph{Say it:} say \emph{le photographe}
  \item \emph{Recall:} say the French for a soldier, then the French for a tradition, then say \emph{le photographe} again
\end{itemize}

\begin{tcolorbox}[breakable,colback=teal!4,colframe=teal!35!black,title={Before you move on}]
What does le photographe mean? (\textquotedblleft{}A photographer\textquotedblright{}.) Say it once more.
\end{tcolorbox}
\section[le journaliste]{le journaliste — a journalist}

\label{lesson:FR-C166-journaliste}



\begin{quote}
Before the new one: say the French for a tradition, then the French for a photographer.
\end{quote}

\subsection*{You'll want to know: le journaliste}
\textbf{le journaliste} — \textquotedblleft{}a journalist\textquotedblright{}.

\begin{grammarlens}[title={What you've built}]
One more word for this chapter.
\end{grammarlens}

\subsection*{Your turn}
\begin{itemize}
\raggedright
  \item \emph{Say it:} \emph{le journaliste}
  \item \emph{Say it:} \emph{le journaliste}, once more
  \item \emph{Say it:} say \emph{le journaliste}
  \item \emph{Recall:} say the French for a tradition, then the French for a photographer, then say \emph{le journaliste} again
\end{itemize}

\begin{tcolorbox}[breakable,colback=teal!4,colframe=teal!35!black,title={Before you move on}]
What does le journaliste mean? (\textquotedblleft{}A journalist\textquotedblright{}.) Say it once more.
\end{tcolorbox}
\section[l'architecte]{l'architecte — an architect}

\label{lesson:FR-C166-architecte}



\begin{quote}
Before the new one: say the French for a photographer, then the French for a journalist.
\end{quote}

\subsection*{You'll want to know: l'architecte}
\textbf{l'architecte} — \textquotedblleft{}an architect\textquotedblright{}.

\begin{grammarlens}[title={What you've built}]
One more word for this chapter.
\end{grammarlens}

\subsection*{Your turn}
\begin{itemize}
\raggedright
  \item \emph{Say it:} \emph{l'architecte}
  \item \emph{Say it:} \emph{l'architecte}, once more
  \item \emph{Say it:} say \emph{l'architecte}
  \item \emph{Recall:} say the French for a photographer, then the French for a journalist, then say \emph{l'architecte} again
\end{itemize}

\begin{tcolorbox}[breakable,colback=teal!4,colframe=teal!35!black,title={Before you move on}]
What does l'architecte mean? (\textquotedblleft{}An architect\textquotedblright{}.) Say it once more.
\end{tcolorbox}
\section[le dentiste]{le dentiste — a dentist}

\label{lesson:FR-C166-dentiste}



\begin{quote}
Before the new one: say the French for a journalist, then the French for an architect.
\end{quote}

\subsection*{You'll want to know: le dentiste}
\textbf{le dentiste} — \textquotedblleft{}a dentist\textquotedblright{}.

\begin{grammarlens}[title={What you've built}]
One more word for this chapter.
\end{grammarlens}

\subsection*{Your turn}
\begin{itemize}
\raggedright
  \item \emph{Say it:} \emph{le dentiste}
  \item \emph{Say it:} \emph{le dentiste}, once more
  \item \emph{Say it:} say \emph{le dentiste}
  \item \emph{Recall:} say the French for a journalist, then the French for an architect, then say \emph{le dentiste} again
\end{itemize}

\begin{tcolorbox}[breakable,colback=teal!4,colframe=teal!35!black,title={Before you move on}]
What does le dentiste mean? (\textquotedblleft{}A dentist\textquotedblright{}.) Say it once more.
\end{tcolorbox}
\section[le pharmacien]{le pharmacien — a pharmacist}

\label{lesson:FR-C166-pharmacien}



\begin{quote}
Before the new one: say the French for an architect, then the French for a dentist.
\end{quote}

\subsection*{You'll want to know: le pharmacien}
\textbf{le pharmacien} — \textquotedblleft{}a pharmacist\textquotedblright{}.

\begin{grammarlens}[title={What you've built}]
One more word for this chapter.
\end{grammarlens}

\subsection*{Your turn}
\begin{itemize}
\raggedright
  \item \emph{Say it:} \emph{le pharmacien}
  \item \emph{Say it:} \emph{le pharmacien}, once more
  \item \emph{Say it:} say \emph{le pharmacien}
  \item \emph{Recall:} say the French for an architect, then the French for a dentist, then say \emph{le pharmacien} again
\end{itemize}

\begin{tcolorbox}[breakable,colback=teal!4,colframe=teal!35!black,title={Before you move on}]
What does le pharmacien mean? (\textquotedblleft{}A pharmacist\textquotedblright{}.) Say it once more.
\end{tcolorbox}
